%!TEX root = main.tex

%%%%%%%%%%%%%%%%%%%%%%%%%%%%%%%%%%%%%%%%%%%%%%%%%%%%%%%%%%%%%
\subsection{The \GNC Algorithm with GM and TLS Costs}
% \subsection{The GNC Algorithm with GM and TLS Costs}
\label{sec:GNC-GM-TLS}

 
%satisfied for a large class of robust cost functions: 
% many penalty functions 
Here we tailor the \GNC algorithm to two cost functions, the 
Geman McClure and the Truncated Least Squares costs, provide expressions for 
the penalty term $\Phi_{\rho_\mu}$, and discuss how to solve the weight update step. 
The proofs of the following propositions are given in the supplementary material~\cite{Yang20tr-GNC}. 
% So far our development has been conceptual and general for many robust cost functions. We now instantiate the GNC algorithm by providing two examples to illustrate the graduated non-convexity governed by control parameter $\mu$, the corresponding outlier processes, and the dual weight update rules for Proposition~\ref{prop:primalDualOutlierProcess}.

% \marginpar{proof or cite}
\begin{proposition}[GNC-Geman McClure (\GNCGM)] \label{prop:GNCGM}
Consider the Geman-McClure function and its \GNC surrogate with control parameter $\mu$, as introduced in Example~\ref{example:GM}. Then, the minimization of the surrogate function $\rho_\mu(\cdot)$ is equivalent to  the 
outlier process with penalty term chosen as:
\bea 
\label{eq:outlierProcessGNCGM}
\Phi_{\rho_\mu}(w_i) = \mu \barcsq (\sqrt{w_i} - 1)^2.
\eea
% \begin{example}[GNC-Geman McClure (\GNCGM)] The Geman-McClure function with control parameter $\mu$ is:
% \bea \label{eq:formulationGNCGM}
% \rho_{\mu}(x) = \frac{\mu x^2}{ \mu + x^2}.
% \eea
% Decreasing $\mu$ increases the non-convexity (see Fig.~\ref{fig:GNCplots}(a)). The outlier process for $\rho_{\mu}(x)$ is:
% \bea \label{eq:outlierProcessGNCGM}
% \Phi(w_i) = \mu(\sqrt{w_i} - 1)^2
% \eea
Moreover, defining the residual $\hatr_i^2 \doteq r^2(\vy_i, \vxx\at{t})$, 
 the weight update at iteration $t$ can be solved in closed form as:
% and the dual weight update rule is:
\bea 
\label{eq:dualWeightUpdateGNCGM}
w_i\at{t} = \left( \frac{\mu \barcsq }{\hatr_i^2 + \mu \barcsq} \right)^2\!.
\eea
\end{proposition}

% \LC{comment on stopping conditions}


% \marginpar{proof or cite}
\begin{proposition}[GNC-Truncated Least Squares (\GNCTLS)] \label{prop:GNCTLS}
Consider the truncated least squares function and its \GNC surrogate with control parameter $\mu$, as introduced in Example~\ref{example:TLS}. Then, the minimization of the surrogate function $\rho_\mu(\cdot)$ is equivalent to  the 
outlier process with penalty term:
% The truncated least squares function with control parameter $\mu$ is:
% \bea
% \hspace{-2mm} \rho(x) = \begin{cases}
% \scriptstyle{ x^2 } &  \scriptstyle{ x^2 \in \left[0, \frac{\mu \barcsq}{\mu + 1} \right] } \\
%   \scriptstyle{ 2\barc | x | \sqrt{\mu(\mu+1)} - \mu (\barcsq + x^2)  } &  \scriptstyle{ x^2 \in \left[\frac{\mu \barcsq}{\mu + 1} , \frac{(\mu+1)\barcsq}{\mu} \right] }\\
%  \scriptstyle{ \barcsq } &  \scriptstyle{ x^2 \in \left[\frac{(\mu+1)\barcsq}{\mu}, +\infty \right) }
% \end{cases}.
% \eea
% Increasing $\mu$ increases the non-convexity (see Fig.~\ref{fig:GNCplots}(b)). The outlier process for $\rho_{\mu}(x)$ is:
\bea \label{eq:outlierProcessGNCTLS}
\Phi_{\rho_\mu}(w_i) = \frac{\mu (1-w_i)}{ \mu + w_i} \barcsq.
\eea
%
Moreover, defining the residual $\hatr_i^2 \doteq r^2(\vy_i, \vxx\at{t})$, 
 the weight update at iteration $t$ can be solved in closed form as:
 %
\bea \label{eq:dualWeightUpdateGNCTLS}
w_i\at{t} = \begin{cases}
0 & \text{ if } \hatr_i^2 \in \left[ \frac{\mu+1}{\mu}\barcsq, +\infty \right] \\
\frac{\barc}{\hatr_i}\sqrt{\mu(\mu+1)} - \mu & \text{ if }  \hatr_i^2 \in \left[ \frac{\mu}{\mu+1}\barcsq ,\frac{\mu+1}{\mu}\barcsq \right] \\
1 &\text{ if }  \hatr_i^2 \in \left[ 0,\frac{\mu}{\mu+1}\barcsq \right].
\end{cases}
\eea
\end{proposition}

\begin{remark}[Implementation Details] %[Initial $\mu$ and Stopping Conditions] 
For \GNCGM, we start with a convex surrogate ($\mu\kern-0.5em\rightarrow\kern-0.5em\infty$) and decrease $\mu$ till we recover the original cost ($\mu\!\rightarrow \!1$). In practice, 
calling $r_{\max}^2 \doteq \max_i( r^2(\vy_i, \vxx^{(0)}))$ the maximum residual after the first variable update,  we initialize $\mu \!= \!2r^2_{\max} / \barcsq$,
 update $\mu \leftarrow \mu / 1.4$ at each outer iteration, 
 and stop when $\mu$ decreases below $1$.
For \GNCTLS, we start with a convex surrogate ($\mu\rightarrow 0$) and increase the $\mu$ till we recover the original cost ($\mu\kern-0.2em\rightarrow\kern-0.2em\infty$). In practice, 
we initialize  $\mu \!= \!\barcsq / ( 2 r^2_{\max} - \barcsq )$,
 update $\mu\!\leftarrow \!1.4 \mu$ at each outer iteration, 
 and stop when the sum of the weighted residuals $\sumAllPointsi w_i \hatr_i^2$ converges.  
 For each outer iteration we perform a single variable and weight update.\revise{~At the first iteration, all weights are set to 1 ($w_i^{(0)}=1, i=1,\dots,N$).} For both robust functions, we set the parameter $\barc$ to the maximum {error expected for the inliers, see Remarks 1-2 in~\cite{Yang19iccv-QUASAR}.}
% The control parameter $\mu$ 
% Define $r_{\max}^2 \doteq \max_i( r^2(\vy_i, \vxx^{(0)}) ) $ as the maximum residual of the first outer iteration, we initialize the initial $\mu$ for \GNCGM as $\mu_0 = r^2_{\max} / e^2$, and the initial $\mu$ for \GNCTLS as $\mu_0 = 1 / ( 2 r^2_{\max} / \barcsq - 1 )$. We stop \GNCGM when $\mu$ decreases below $1$, and we stop \GNCTLS when the sum of the weighted residuals $\sumAllPointsi w_i \hatr_i^2$ converges. 
%\HY{The next section will show the convergence for \GNCTLS.}
\end{remark}
% \red{Proof for the dual weight update rules are presented in the Appendix.}

